% !TeX root = ../main.tex
\section{Vectorization Verification}\label{sec:vectorization}
\rubric{Vectorization verification}{15}{Clear evidence that SIMD instructions are used, with a correct
reading of vector width, FMA and masking.}

\begin{hypothesis}
  \todo{Which loops you expect each variant to vectorize, at what width, and which you expect the
  compiler to reject (and why).}
\end{hypothesis}

\Cref{tab:vec-summary} counts, for each kernel function, the instructions that operate on vector
registers or use packed operations (\code{make report}).

\begin{table}[tbp]
  \centering\small
  \caption{SIMD instructions per kernel function and build variant.}\label{tab:vec-summary}
  \IfDataExists{vectorization}{%
    \pgfplotstabletypeset[
      col sep=comma,
      columns={variant,function,insns,simd_insns,widest_register},
      columns/variant/.style={column name=Variant, string type},
      columns/function/.style={column name=Function, string type, string replace*={_}{\_}},
      columns/insns/.style={column name=Instructions},
      columns/simd_insns/.style={column name=SIMD instructions},
      columns/widest_register/.style={column name=Widest register, string type},
      every head row/.style={before row=\toprule, after row=\midrule},
      every last row/.style={after row=\bottomrule},
    ]{data/vectorization.csv}%
  }
\end{table}

\subsection{Compiler report}
\begin{lstlisting}[style=excerpt]
TODO: 2-5 relevant lines from data/raw/vectorization/<run>/report_<variant>.txt
\end{lstlisting}

\subsection{Disassembly}
\begin{lstlisting}[style=excerpt]
TODO: 5-10 lines of a vectorized loop body, e.g. disasm/native/saxpy_f32.s
\end{lstlisting}

\todo{Interpret the evidence: register width (xmm/ymm/zmm) and lanes per element type; FMA
(\code{vfmadd*}); masked tails (k-mask registers) versus a scalar epilogue; loops that did not
vectorize and the compiler's reason.}
